\documentclass[a4paper, serif, 10pt]{moderncv}

%% moderncv
% style and color
\moderncvstyle{banking}
\moderncvcolor{black}

%% page layout/margins
\usepackage[top=2.5cm, bottom=2.5cm, left=1.5cm, right=1.5cm]{geometry}

\nopagenumbers{}

%% Babel config for Dutch language
% ref:
% - https://latex3.github.io/babel/guides/locale-dutch.html
\usepackage[dutch]{babel}

%% fontspec
\usepackage{fontspec}

\IfFontExistsTF{Linux Libertine}{
  \setmainfont[Ligatures={TeX, Common}, Numbers=OldStyle]{Linux Libertine}
}{}
\IfFontExistsTF{Linux Libertine O}{
  \setmainfont[Ligatures={TeX, Common}, Numbers=OldStyle]{Linux Libertine O}
}{}

%% CTeX
% CJK support, used to show CJK characters in the resume
%
% - fontset=none: disable builtin fontset but instead set the CJK font manually
% - heading=false: disable ctex heading
% - punct=kaiming: use kaiming punctuations styles for CJK
% - scheme=plain: use plain scheme, do not override \normalsize font size
% - space=auto: space settings for CJK characters
%
% ref:
% - http://ctan.mirrorcatalogs.com/language/chinese/ctex/ctex.pdf
\usepackage[UTF8, heading=false, punct=kaiming, scheme=plain, space=auto]{ctex}

\IfFontExistsTF{Noto Serif CJK SC}{
  \setCJKmainfont{Noto Serif CJK SC}
}{}
\IfFontExistsTF{Noto Sans CJK SC}{
  \setCJKsansfont{Noto Sans CJK SC}
}{}

%% line spacing
\usepackage{setspace}
\setstretch{1.125}

%% URL styling - use normal text instead of monospace
\urlstyle{same}

% Auto-underline all links
% Defer until \begin{document} so hyperref is already loaded
\AtBeginDocument{%
  \let\oldhref\href
  \renewcommand{\href}[2]{\oldhref{#1}{\underline{#2}}}%
}

%% Patch \httplink and \httpslink to handle full URLs
% The original moderncv macros always prepend http:// or https:// to URLs.
% This causes issues when the URL already has a protocol (e.g., https://example.com)
% resulting in malformed URLs like https://https://example.com.
% This patch checks if the URL already contains "://" and uses it directly if so.
% Note: We use \str_set:Nx (x-expansion) to expand macros like \@homepage first.
\ExplSyntaxOn
\str_new:N \g_yamlresume_url_str

\RenewDocumentCommand{\httpslink}{O{}m}{
  \str_set:Nx \g_yamlresume_url_str {#2}
  \str_if_in:NnTF \g_yamlresume_url_str {://}
    { \url{#2} }
    { \url{https://#2} }
}

\RenewDocumentCommand{\httplink}{O{}m}{
  \str_set:Nx \g_yamlresume_url_str {#2}
  \str_if_in:NnTF \g_yamlresume_url_str {://}
    { \url{#2} }
    { \url{http://#2} }
}
\ExplSyntaxOff

\name{Thomas de Vries}{}
\title{DevOps Engineer | Cloudinfrastructuur \& Automatisering}
\phone[mobile]{+31 20 123 4567}
\email{thomas.devries@example.com}
\homepage{https://thomasdevries.dev}

\address{Keizersgracht 123, Amsterdam, Noord-Holland, Nederland, 1015 CJ}{}{}

\extrainfo{{\small \faLinkedin}\ \href{https://www.linkedin.com/in/thomasdevries}{@thomasdevries} {} {} {} • {} {} {} 
{\small \faGithub}\ \href{https://github.com/thomasdevries}{@thomasdevries} {} {} {} • {} {} {} 
{\small \faStackOverflow}\ \href{https://stackoverflow.com/users/12345678/thomasdevries}{@thomasdevries}}

\begin{document}

\maketitle

\section{Basis Informatie}

\cvline{}{\begin{itemize}
\item Ervaren DevOps Engineer met meer dan 8 jaar ervaring in het automatiseren van cloudinfrastructuur en het bouwen van betrouwbare CI/CD-pipelines.
\item Diepgaande kennis van AWS, Azure, Kubernetes, Terraform en GitOps.
\item Gepassioneerd door platform engineering, observability en het verhogen van de ontwikkelsnelheid van teams.
\item Sterke communicatieve vaardigheden en ervaring met het begeleiden van multidisciplinaire teams.
\end{itemize}
Ik werk graag aan de snijlijn van infrastructuur en softwareontwikkeling, waarbij ik processen continu verbeter en operationele kosten verlaag zonder in te leveren op betrouwbaarheid.}
\section{Opleidingen}

\cventry{sep 2016–jul 2018}
        {Master, Computerwetenschappen, Score: 8.5}
        {Technische Universiteit Delft}
        {\url{https://www.tudelft.nl/}}
        {}
        {\begin{itemize}
\item Afstudeerrichting Cloud \& Distributed Systems
\item Thesis over automatische schaalvergroting van containerplatforms in hybride clouds
\item Cum laude afgestudeerd
\end{itemize}
\textbf{Cursussen}: Gedistribueerde systemen, Cloud computing, Netwerktechnologie, Operationele systemen, Softwarearchitectuur, Beveiliging van informatiesystemen}

\cventry{sep 2012–jun 2016}
        {Bachelor, Informatica, Score: 8.0}
        {Hogeschool van Amsterdam}
        {\url{https://www.hva.nl/}}
        {}
        {\begin{itemize}
\item Minor in Cloud \& Virtualisatie
\item Stage bij een hostingprovider voor geautomatiseerde deployments
\end{itemize}
\textbf{Cursussen}: Programmeren in Python, Linux systeembeheer, Databases, Netwerken, IT-beveiliging}
\section{Werkervaring}

\cventry{jan 2022–Heden}
        {Senior DevOps Engineer}
        {Rabobank}
        {\url{https://www.rabobank.nl/}}
        {}
        {\begin{itemize}
\item Verantwoordelijk voor het ontwerpen en beheren van een multi-cloud platform op AWS en Azure.
\item Geautomatiseerde CI/CD-pipelines met GitLab CI en Argo CD, waardoor releases van dagen naar minuten gingen.
\item Infrastructure as Code met Terraform en Ansible voor meer dan 200 applicaties.
\item Implementatie van observability met Prometheus, Grafana en Loki.
\item Begeleiden van een team van 6 DevOps engineers en bevorderen van kennisdeling.
\end{itemize}
\textbf{Trefwoorden}: AWS, Azure, Terraform, Kubernetes, GitLab CI, Observability}

\cventry{aug 2018–dec 2021}
        {DevOps Engineer}
        {KPN}
        {\url{https://www.kpn.com/}}
        {}
        {\begin{itemize}
\item Beheerde Kubernetes-clusters en ontwikkelde Helm-charts voor interne microservices.
\item Bouwde CI/CD-pipelines met Jenkins en later GitHub Actions.
\item Automatiseerde serverprovisioning met Ansible en Packer.
\item Verhoogde beschikbaarheid van diensten naar 99,95\% door monitoring en alerting.
\item Samenwerking met ontwikkelteams om security en compliance te integreren in de pipeline.
\end{itemize}
\textbf{Trefwoorden}: Kubernetes, Jenkins, Ansible, Packer, Prometheus}

\cventry{sep 2015–jul 2018}
        {Systeembeheerder / DevOps}
        {Info Support}
        {\url{https://www.infosupport.com/}}
        {}
        {\begin{itemize}
\item Beheer van Linux- en Windows-servers voor zakelijke klanten.
\item Introductie van geautomatiseerde configuratie met Puppet.
\item Migratie van on-premises workloads naar AWS.
\item Opzetten van basis monitoring met Nagios en later Zabbix.
\end{itemize}
\textbf{Trefwoorden}: Linux, Windows Server, Puppet, AWS, Nagios}
\section{Talen}

\cvline{Nederlands}{Moedertaal of tweetalige vaardigheid \hfill \textbf{Trefwoorden}: Moedertaal}
\cvline{Engels}{Volledige professionele vaardigheid \hfill \textbf{Trefwoorden}: Professioneel werkzaam, TOEIC 950}
\cvline{Duits}{Beperkte werkvaardigheid \hfill \textbf{Trefwoorden}: Basiskennis, A2}
\section{Vaardigheden}

\cvline{Cloudplatformen}{Expert \hfill \textbf{Trefwoorden}: AWS, Azure, Google Cloud}
\cvline{Containers \& Orchestratie}{Expert \hfill \textbf{Trefwoorden}: Kubernetes, Docker, Helm, OpenShift}
\cvline{Infrastructure as Code}{Geavanceerd \hfill \textbf{Trefwoorden}: Terraform, Ansible, Pulumi, CloudFormation}
\cvline{CI/CD \& Automatisering}{Expert \hfill \textbf{Trefwoorden}: GitLab CI, Jenkins, GitHub Actions, Argo CD}
\cvline{Monitoring \& Observability}{Geavanceerd \hfill \textbf{Trefwoorden}: Prometheus, Grafana, Loki, ELK Stack}
\cvline{Programmeertalen \& Scripting}{Geavanceerd \hfill \textbf{Trefwoorden}: Python, Bash, Go, PowerShell}
\section{Onderscheidingen}

\cventry{nov 2023}
        {Rabobank}
        {DevOps Team van het Jaar}
        {}
        {}
        {\begin{itemize}
\item Toegekend voor het succesvol migreren van kritische applicaties naar een geautomatiseerd multi-cloud platform.
\end{itemize}}

\cventry{jul 2018}
        {Technische Universiteit Delft}
        {Beste Afstudeerproject}
        {}
        {}
        {\begin{itemize}
\item Prijs voor onderzoek naar automatische schaalvergroting van containerplatforms.
\end{itemize}}
\section{Certificaten}

\cventry{mrt 2023}
        {AWS}
        {AWS Certified DevOps Engineer - Professional}
        {\url{https://aws.amazon.com/certification/certified-devops-engineer-professional/}}
        {}
        {}

\cventry{jun 2021}
        {CNCF}
        {Certified Kubernetes Administrator (CKA)}
        {\url{https://www.cncf.io/certification/cka/}}
        {}
        {}

\cventry{sep 2022}
        {HashiCorp}
        {HashiCorp Certified: Terraform Associate}
        {\url{https://www.hashicorp.com/certification/terraform-associate}}
        {}
        {}
\section{Publicaties}

\cventry{mei 2023}
        {Zero Trust in cloud-native omgevingen}
        {Dutch DevOps Journal}
        {\url{https://www.dutchdevopsjournal.nl/zero-trust-cloud-native}}
        {}
        {\begin{itemize}
\item Artikel over het implementeren van een Zero Trust-architectuur in Kubernetes-omgevingen.
\item Gepubliceerd in het vaktijdschrift Dutch DevOps Journal.
\end{itemize}}
\section{Projecten}

\cventry{mrt 2022–okt 2023}
        {Inner source platform voor self-service cloudinfrastructuur}
        {Platform Engine}
        {\url{https://github.com/thomasdevries/platform-engine}}
        {}
        {\begin{itemize}
\item Ontwikkelde een self-service platform waarmee ontwikkelteams binnen 10 minuten een productieklare omgeving konden aanvragen.
\item Gebruikt Terraform, Kubernetes en Crossplane voor infrastructure provisioning.
\item Verlaagde de doorlooptijd van omgevingen met 70\%.
\end{itemize}
\textbf{Trefwoorden}: Platform Engineering, Self-service, Crossplane, Kubernetes}

\cventry{jan 2021–dec 2021}
        {Migratie van push-based naar pull-based deployments}
        {GitOps-migratie}
        {\url{https://github.com/thomasdevries/gitops-migratie}}
        {}
        {\begin{itemize}
\item Introduceerde Argo CD voor declaratieve deployments.
\item Verhoogde de deployment-frequentie en verlaagde de foutmarge.
\end{itemize}
\textbf{Trefwoorden}: GitOps, Argo CD, Kubernetes}
\section{Interesses}

\cvline{Sport}{Hardlopen, Wielrennen, Zwemmen}
\cvline{Muziek}{Gitaar, Piano}
\cvline{Open source}{Kubernetes, Terraform, CNCF}
\section{Vrijwilligerswerk}

\cventry{jan 2019–Heden}
        {Community Organizer}
        {DevOps Nederland}
        {\url{https://www.devopsnederland.nl/}}
        {}
        {\begin{itemize}
\item Organiseer maandelijkse meetups en workshops over DevOps, cloud en Kubernetes.
\item Verantwoordelijk voor sprekers, locatie en community-opbouw.
\item Groei van de community van 300 naar 1200 leden.
\end{itemize}}
    
\end{document}